\scanpage{92}
l'homomorphisme canonique
\[
\tag{45}\mathcal C'\longrightarrow\widetilde{\mathcal C}^{\unsure{\prime|*}}.
\]
\inserted{vrai revêtement de $\mathcal C$, dont les sommets sont les $\Pi_0$ des fibres de $\mathcal C'\mathbin{-}\widetilde{\mathcal C}^{\unsure{*|\prime}}$}
\deleted{\unsure{soit}} factorisant $\mathcal C'\mathbin{-}\mathcal C$ est injectif, de sorte qu'on peut identifier $\mathcal C'$ à un sous-groupoïde plein de $\widetilde{\mathcal C}^{*}$, équivalent à $\widetilde{\mathcal C}^{*}$ par l'inclusion, où $\widetilde{\mathcal C}$ est un rev. fini de $\mathcal C$\unsure{.}

Soit $x'_0$ un sommet de $\mathcal C'$, $x_0$ son image dans $\widetilde{\mathcal C}$. Soit $\Theta$ un automorphisme de $\mathcal C$ qui fixe $x_0$, \deleted{alors s'}il \deleted{\unsure{y en a un de}\unsure{.}} \inserted{revient au} $=$ de se donner un autom. de \deleted{$\mathcal C$} $\Pi_1(\mathcal C,x_0)\overset{\mathrm{def}}{=}\Pi_0$ (\unsure{abrég.}), et pour tout \inserted{sommet} $\deleted{\unsure{x}}x_i\in\operatorname{Ob}\mathcal C$, $x_i\ne x_0$, un automorphisme du $\Pi_1(x_0)$-torseur à droite $H_{x_i,x_0}=\operatorname{Hom}_{\mathcal C}(x_0,x_i)$, compatible avec l'opération de $\Theta$ sur $\Pi_0$. Ceci posé, pour que cette action se remonte en une action de $\widetilde{\mathcal C}$ qui fixe $x'_0$, il f et s. que le sous-groupe $\Pi'_0$ de $\Pi_0$, image de \deleted{$\unsure{\Pi_1(}$}
\[
\Pi_1(\widetilde{\mathcal C},x'_0)\hookrightarrow\Pi_1(\mathcal C,x_0)=\Pi_0,
\]
soit \uline{stable} par $\Theta$ --- l'action de $\Theta$ sur les fibres \inserted{$\widetilde{\mathcal C}_{x_i}$} de $\widetilde{\mathcal C}/\mathcal C$ est alors donnée par la formule
\[
\tag{45}\operatorname{Ob}\widetilde{\mathcal C}_{x_i}=H_{x_i,x_0}\wedge_{\Pi_0}\underbrace{\Pi_0/\Pi'_0}_{\substack{\text{\unsure{espace|esp.} \unsure{hom|hom.}}\\\text{des classes}\\\text{à droite, sur lequel}\\\Pi_0\ \text{opère à \unsure{g|g.}}}}
\]
par transport de structure : partir de son
\margin{91}
